\documentclass{report}
\usepackage{amsmath,amssymb}
\setlength{\parindent}{0mm}
\setlength{\parskip}{1em}
\begin{document}
\begin{center}
\rule{6in}{1pt} \
{\large Jodi Mead \\
{\bf Regularization of nonlinear inverse problems}}

Mathematics Department \\ Boise State University \\ 1910 University Dr \\ Boise \\ ID 83725-1555
\\
{\tt jmead@boisestate.edu}\\
Chad Hammerquist\end{center}

Tikhonov regularization is commonly used to solve nonlinear ill-posed
inverse problems. The convergence of these methods is verified as the
noise in the data goes to zero. However, in most applications the data
noise does not go zero, and regularization parameters based on noise
estimates are found. For example, the discrepancy principle can be viewed
as applying a $\chi^2$ test to the data residual in order to determine
the regularization parameter. However, when the number of parameters is
greater than or equal to the number of data it is not possible to apply
the $\chi^2$ test because the degrees of freedom is negative or zero. We
suggest applying the $\chi^2$ test to the regularized residual because
the degrees of freedom is equal to the number of data, and we call this
approach the {\it $\chi^2$ method}.

In this talk we describe how the $\chi^2$ method can be applied to
nonlinear problems. We will show both analytically and numerically that
each iterate of of the Gauss-Newton and Levenburg-Marquart methods follow
a $\chi^2$ distribution. This property can be used to find regularization
parameters as is done with the discrepancy principle in Occam's method.
However, our approach differs in that we used the regularized residual
rather than the data residual, and can estimate a different
regularization parameter at each iterate.


\end{document}
